Let $t\ge1$ and $u\in\RE\setminus\{\pm1\}$ with $u\equiv1\pmod{2^t\mathcal O_n}$. By Lemma~\ref{lem:A}, $u\notin B_{n-1}$; since $B_n/\mathbb Q$ is cyclic of $2$-power degree, its subfields are the $B_k$, so $\mathbb Q(u)=B_n$ and $u$ has degree $N=2^n$. From $u-1\in2^t\mathcal O_n$ and $u+1=(u-1)+2\in2\mathcal O_n$ we get $u^2-1\in2^{t+1}\mathcal O_n$, so the ideal $\mathfrak M=2^{t+1}\mathcal O_n$ contains $u^2-1$ and has absolute norm $C=2^{(t+1)N}$, $C^{1/N}=2^{t+1}$. Schinzel's inequality in the form quoted as \cite{MO3} Thm.~5.1 (``Let $\varepsilon$ be a totally real unit different from $\pm1$ and $\mathfrak M$ an ideal of $\mathbb Q(\varepsilon)$ containing $\varepsilon^2-1$; then $M(\varepsilon)\ge\big((C^{1/d}+\sqrt{C^{2/d}+4})/2\big)^{d/2}$, $C$ the absolute norm of $\mathfrak M$, $d$ the degree'') gives, with $d=N$, $M(u)\ge\big((2^{t+1}+\sqrt{4^{t+1}+4})/2\big)^{N/2}=\big(2^t+\sqrt{4^t+1}\big)^{N/2}=\exp\big(\tfrac N2\operatorname{arcsinh}(2^t)\big)$, using $\operatorname{arcsinh}x=\log(x+\sqrt{x^2+1})$. Finally, for a unit the conjugates satisfy $\sum_i\log|u_i|=0$, so $\log M(u)=\sum_{|u_i|>1}\log|u_i|=\tfrac12\sum_i|\log|u_i||\le\tfrac12\sqrt N\,\hgt(u)$ by Cauchy--Schwarz (this is \cite{MO16} eq.~(2.4), proof of their Lemma~2.3). Hence $\hgt(u)\ge\frac{2}{\sqrt N}\cdot\frac N2\operatorname{arcsinh}(2^t)=\sqrt{2^n}\operatorname{arcsinh}(2^t)=L_{n,t}$. For $t=1$ this is $\sqrt{2^n}\log(2+\sqrt5)=L_n$, because $\operatorname{arcsinh}2=\log(2+\sqrt5)$. (This floor is proved here for every $t\ge1$, so Theorem~\ref{thm:S0} keeps only hypothesis (i), which remains open for $t\ge5$.)
